\documentclass[10pt,a4paper]{amsart}
\usepackage{amsmath,amssymb,amsthm,graphicx}
\usepackage{verbatim}
\newenvironment{reference}{\comment}{\endcomment}
%\newenvironment{reference}{}{}
\newenvironment{slogan}{\comment}{\endcomment}
\newenvironment{history}{\comment}{\endcomment}

% For commutative diagrams we use Xy-pic
\usepackage[all]{xy}

% We use 2cell for 2-commutative diagrams.
\xyoption{2cell}
% Theorem environments.
%
\theoremstyle{plain}
\newtheorem{theorem}[subsection]{Theorem}
\newtheorem{proposition}[subsection]{Proposition}
\newtheorem{lemma}[subsection]{Lemma}

\theoremstyle{definition}
\newtheorem{definition}[subsection]{Definition}
\newtheorem{example}[subsection]{Example}
\newtheorem{exercise}[subsection]{Exercise}
\newtheorem{situation}[subsection]{Situation}

\theoremstyle{remark}
\newtheorem{remark}[subsection]{Remark}
\newtheorem{remarks}[subsection]{Remarks}

\numberwithin{equation}{subsection}

% Macros
%
\def\lim{\mathop{\mathrm{lim}}\nolimits}
\def\colim{\mathop{\mathrm{colim}}\nolimits}
\def\Spec{\mathop{\mathrm{Spec}}}
\def\Hom{\mathop{\mathrm{Hom}}\nolimits}
\def\Ext{\mathop{\mathrm{Ext}}\nolimits}
\def\SheafHom{\mathop{\mathcal{H}\!\mathit{om}}\nolimits}
\def\SheafExt{\mathop{\mathcal{E}\!\mathit{xt}}\nolimits}
\def\Sch{\mathit{Sch}}
\def\Mor{\mathop{\mathrm{Mor}}\nolimits}
\def\Ob{\mathop{\mathrm{Ob}}\nolimits}
\def\Sh{\mathop{\mathit{Sh}}\nolimits}
\def\NL{\mathop{N\!L}\nolimits}
\def\CH{\mathop{\mathrm{CH}}\nolimits}
\def\proetale{{pro\text{-}\acute{e}tale}}
\def\etale{{\acute{e}tale}}
\def\QCoh{\mathit{QCoh}}
\def\Ker{\mathop{\mathrm{Ker}}}
\def\Im{\mathop{\mathrm{Im}}}
\def\Coker{\mathop{\mathrm{Coker}}}
\def\Coim{\mathop{\mathrm{Coim}}}

% Boxtimes
%
\DeclareMathSymbol{\boxtimes}{\mathbin}{AMSa}{"02}

%
% Macros for moduli stacks/spaces
%
\def\QCohstack{\mathcal{QC}\!\mathit{oh}}
\def\Cohstack{\mathcal{C}\!\mathit{oh}}
\def\Spacesstack{\mathcal{S}\!\mathit{paces}}
\def\Quotfunctor{\mathrm{Quot}}
\def\Hilbfunctor{\mathrm{Hilb}}
\def\Curvesstack{\mathcal{C}\!\mathit{urves}}
\def\Polarizedstack{\mathcal{P}\!\mathit{olarized}}
\def\Complexesstack{\mathcal{C}\!\mathit{omplexes}}
% \Pic is the operator that assigns to X its picard group, usage \Pic(X)
% \Picardstack_{X/B} denotes the Picard stack of X over B
% \Picardfunctor_{X/B} denotes the Picard functor of X over B
\def\Pic{\mathop{\mathrm{Pic}}\nolimits}
\def\Picardstack{\mathcal{P}\!\mathit{ic}}
\def\Picardfunctor{\mathrm{Pic}}
\def\Deformationcategory{\mathcal{D}\!\mathit{ef}}
\usepackage{cleveref}
\crefname{originaldefinition}{original definition}{original definitions}
\Crefname{originaldefinition}{Original definition}{Original definitions}
\makeatletter
\newlabel{divisors-definition-admissible-blowup}{{080K}{source}{}{sourcepointer-1}{}}
\newlabel{divisors-definition-admissible-blowup@cref}{{[originaldefinition][0][]080K}{[1][0][]source}{}{}{}}
\makeatother
\crefname{originallemma}{original lemma}{original lemmas}
\Crefname{originallemma}{Original lemma}{Original lemmas}
\makeatletter
\newlabel{divisors-lemma-flat-base-change-blowing-up}{{0805}{source}{}{sourcepointer-2}{}}
\newlabel{divisors-lemma-flat-base-change-blowing-up@cref}{{[originallemma][0][]0805}{[1][0][]source}{}{}{}}
\makeatother
\title{Descending admissible blowups of the punctured spectrum along the completion}
\author{The Stacks Project Authors}
\date{stacks-project:resolve.tex:local-snapshot-20261012; GNU FDL 1.2 or later}
\begin{document}
\maketitle
\noindent\textit{Editorial note.} All delivered mathematical text is the author's own, copied from the pinned Stacks Project chapter file resolve.tex; no statement or proof step was rewritten, reordered or removed.\par\medskip
\noindent\textit{Editorial note.} The counted proof is the author's complete argument: the descent of the defining ideal through the comparison of the quotients by the powers of the maximal ideal, and the flat base change identification of the blowing up.\par\medskip
\noindent\textit{Editorial note.} The label lemma-descend-admissible-blowup names the counted statement, and the proof environment immediately following it in the source is delivered verbatim as its proof.\par\medskip
\noindent\textit{Editorial note.} The counted statement is carried by the author's own lemma environment in the chapter source.\par\medskip
\noindent\textit{Editorial note.} Pre-open gate: python operations/agent-authoring/tools/precard\_gate.py --source resolve.tex --label lemma-descend-admissible-blowup --start 3028 --end 3050 --local-id 28114 returned RESULT=OK (L28114 free, occupancy FREE, public interval clear). Pre-assembly tag self-check: canonical tag resolved and not present in reference/public-items.jsonl.\par\medskip
\noindent\textit{Editorial note.} The statement and proof are copied from the pinned chapter file; the cited inputs are kept as external pointers with their canonical Stacks tags.\par\medskip
\noindent\textit{Allowed standard prerequisites (editorial).}\begin{itemize}
\item Definition of an admissible blowup. Divisors, the definition of an admissible blowing up with respect to an open subscheme. Provides the ideal whose vanishing locus is the closed point and whose blowing up is the given morphism.
\item Comparison of the base with its completion modulo powers of the maximal ideal. The identity of the quotient of the completion by the powers of the extended maximal ideal with the corresponding quotient of the base. Descends the ideal defining the blowup over the completion to an ideal of the base containing a power of the maximal ideal.
\item Base change of blowing up along a flat morphism. Divisors, the lemma that blowing up commutes with flat base change. Identifies the given morphism over the completion with the base change of the constructed blowing up over the base.
\end{itemize}\medskip
\ifdefined\hypertarget\hypertarget{sourcepointer-1}{}\fi
\noindent\textit{Source reference.} Original definition 080K: sources/stacks-chow/tags.txt, source lines 10185--10185. Prerequisite pointer invoked at proof line(s) 3041; it names a standard fact of a sibling chapter that the counted argument uses, and every step of the counted proof built on it is retained verbatim in the delivered body.\par\medskip
\ifdefined\hypertarget\hypertarget{sourcepointer-2}{}\fi
\noindent\textit{Source reference.} Original lemma 0805: sources/stacks-chow/tags.txt, source lines 10170--10170. Prerequisite pointer invoked at proof line(s) 3049; it names a standard fact of a sibling chapter that the counted argument uses, and every step of the counted proof built on it is retained verbatim in the delivered body.\par\medskip
\begin{lemma}
\label{lemma-descend-admissible-blowup}
Let $(A, \mathfrak m)$ be a Noetherian local ring with completion $A^\wedge$.
Let $U \subset \Spec(A)$ and $U^\wedge \subset \Spec(A^\wedge)$ be the
punctured spectra. If $Y \to \Spec(A^\wedge)$ is a $U^\wedge$-admissible
blowup, then there exists a $U$-admissible blowup $X \to \Spec(A)$
such that $Y = X \times_{\Spec(A)} \Spec(A^\wedge)$.
\end{lemma}

\begin{proof}
By definition there exists an ideal $J \subset A^\wedge$ such that
$V(J) = \{\mathfrak m A^\wedge\}$ and such that $Y$ is the blowup
of $S^\wedge$ in the closed subscheme defined by $J$, see
Divisors, Definition \ref{divisors-definition-admissible-blowup}.
Since $A^\wedge$ is Noetherian this implies
$\mathfrak m^n A^\wedge \subset J$ for some $n$.
Since $A^\wedge/\mathfrak m^n A^\wedge = A/\mathfrak m^n$
we find an ideal $\mathfrak m^n \subset I \subset A$
such that $J = I A^\wedge$. Let $X \to S$ be the blowup in $I$.
Since $A \to A^\wedge$ is flat
we conclude that the base change of $X$ is $Y$ by
Divisors, Lemma \ref{divisors-lemma-flat-base-change-blowing-up}.
\end{proof}

\end{document}
